\documentclass[leqno, a4paper, 12pt]{jsarticle}
\usepackage{amssymb}
\usepackage{amsthm}
\usepackage{amsmath}
\usepackage{comment}
\usepackage{url}

\usepackage{enumitem}
%\usepackage{showkeys}


%定理環境 thmはあまり使わずpropをなるべく使う。
%主張は基本的に証明の中で使い、そのため番号を付けない。
%注意環境を付けてみた。
\newtheorem{thm}{定理}
\newtheorem{prop}[thm]{命題}
\newtheorem{cor}[thm]{系}
\newtheorem{lem}[thm]{補題}
\newtheorem{df}[thm]{定義}
\newtheorem{exam}[thm]{例}
\newtheorem{fact}[thm]{事実}
\newtheorem*{claim}{主張}
\newtheorem*{cau}{注意}
\newtheorem*{rmk}{注意}
\renewcommand{\proofname}{証明}
%矢印たち。
%\toは\rightarrowと同じ。\getsは\leftarrowと同じ。同値は\iff
\newcommand{\To}{\Longrightarrow}
\newcommand{\Gets}{\LongLeftarrow}
%黒板太字
\newcommand{\nn}{\mathbb{N}}
\newcommand{\zz}{\mathbb{Z}}
\newcommand{\qq}{\mathbb{Q}}
\newcommand{\rr}{\mathbb{R}}
\newcommand{\cc}{\mathbb{C}}
%無限大
\newcommand{\mgn}{\infty}
%差集合は\setminusで統一する。等号の否定は\neq
%真部分集合は\subsetneqq　偏微分は\partial
%ディスプレイスタイル。\dfracでディスプレイ分数
\newcommand{\dpst}{\displaystyle}

\title{論文メモ
\large{\\--- 複素関数論基礎論発展編 ---}
}
\author{ナンブ　キトラ}
\date{\today}
\begin{document}
\maketitle

本棚を整理したいので，
溜まっている論文のメモを書いて未来への手紙とすることにする．



この論文の束は複素関数論に関するものである．

\section{序論}

\subsection{コーシー・リーマン方程式}
複素関数が微分可能であれば
コーシー・リーマンの方程式
を満たすことがすぐにわかるが，
複素数値関数の実部や虚部がコーシー・リーマンの方程式を
満たしていたとして，もとの関数それ自体が
複素関数として微分可能かどうかは全然自明ではない．
なぜならコーシー・リーマンの方程式に出てくる
偏微分は$x$軸方向と$y$軸方向のみ
考えているので，このように考えると，
複素平面に存在する全ての方向からの微分が存在するかは
とても微妙である．
よって何かしらの条件を
課した上でコーシー・リーマンの方程式から
複素関数の微分可能性を復元する研究がある．
論文\cite{MR0470179}
によれば，
グルサの結果として
$f=u+v\sqrt{-1}$
がコーシー・リーマンの方程式を満たし
なおかつ$f$が連続で$u$と
$v$の偏微分も連続ならば$f$
は解析的である．
また
Looman--Menchoff
の結果も紹介されている．
それは$f=u+v\sqrt{-1}$
について，$u$と
$v$の偏微分が各点で存在し
なおかつ各点でコーシー・リーマンの方程式を満たし，
さらに$f$が連続であれば
$f$は解析的であることを述べている．
偏微分の連続性に全く言及していないことに注意しよう．
論文\cite{MR0470179}ではこのような結果が色々述べられているが，
Looman--Menchoffはやっぱり一番いいのではないだろうか．
論文
\cite{MR1057270}
では
Looman--Menchoff
の結果の拡張として
$f$に測度論的な仮定，
例えばある$L^{p}$に属するとか，
一次元ハウスドルフ測度有限な閉集合の可算和集合を除いて
偏微分が存在するとかの仮定のもとで
関数の解析性を証明している．
解析関数というとても綺麗な対象を
測度論だとかそういう混沌とした仮定から導くのはとても
面白い分野だと思う．



\subsection{自然境界}

昔の「日本数学物理学会誌」
は外国の論文の翻訳とかを載せていたようだ．
論文
\cite{1927225}
もモーデルによる自然境界をもつ冪級数の
論文の翻訳である．
複素関数$f$の自然境界とは
$f$が解析接続できないような境界になっている
集合のことだったと思う．
$l(n)$が「いい感じ」に発散する自然数列だとすると
$\sum_{n=1}^{\infty}b_{n}x^{l(n)}$
(間隙級数)は$|z|=1$を自然境界をもつというのがアダマールの定理だそうで，
この論文
\cite{1927225}
ではその簡単な証明を与えているようだ．
何でこれをわざわざ印刷して持っていたのかは知らない．



\subsection{リウビルの定理}
リウビルの定理というのは
$\cc$上の解析的複素関数が
有界ならば定数になるというものである．
調和関数の同じ性質のことも
リウビルの定理と呼ばれていると思う．
論文
\cite{MR0259149}
では
リウビルの定理の
証明を与えているのだが，
この論文9行しかない．タイトルや注釈などを含めても
13行ほどしかなく，印刷しても紙の半分程度しか占めないであろう．そういう珍しさから印刷したんだと思う．

論文
\cite{MR0108809}
はリウビルはリウビルでも違う定理の話をしている．
等角写像が$C^{1}$になるみたいなこともリウビルの定理というらしい．それでこれを高次元に拡張して
$C^{\sigma}$
級のリーマン計量を持つリーマン多様体の間の$C^{1}$等長写像は
$C^{1+\sigma}$
ということを証明しているようだ





\subsection{リーマンの写像定理}
複素平面
$\cc$
の単連結な
領域は全て正則同型であるというのが
リーマンの
写像定理である．
この話はリーマン
面の分類とかそういう話にも関わってくると思うがそれはさておき
論文を紹介していく．
歴史については
\cite{MR0323996}
を見ると良いと思う．
証明としてはモンテルの定理を
使って謎の正規族を構成して写像を構成すると思う．
このモンテルの定理というのは後から出てくる．
リーマンの写像定理を
境界の点まで拡張しようという話がありそれは
カラテオドリの定理と呼ばれている．
曰く，単位開円盤から伸びる双正則写像が
境界まで拡張できるための必要十分条件は
行き先の境界がジョルダン閉曲線であることらしい．
この定理を頑張って証明しようとして
境界要素とかprime endsとかの概念が発明されたと思う．
以下の論文を参考のこと．
\cite{qiu2013riemann}
\cite{MR0100090}
論文
\cite{MR0614728}
の
prime endsはなんだかいい概念だった気がする．
\cite{MR0145063}も参照のこと．
\cite{MR0272994}．
現代的には代数的トポロジーが発達しているからそのような道具を使ってこれらの話を書き直すようなことができるのだろうか？


\section{本論}
ここからどんどんいい話になる．
\subsection{正規族，モンテルの定理，ブロッホの原理}

とりあえず
\cite{MR0379852}
と
\cite{MR2241035}
を読もう．とても良い論文．
まず正規族というのを定義する．
領域$D\subset \cc$上で定義された
複素関数からなる族
$\mathcal{F}$
が
\textbf{正規}であるとは，
流儀の違いはあるだろうが(めんどくさいので私はコンパクト開位相で相対コンパクトと定義したいのだが)，
\cite{MR0379852}を参考にすると，
これに属する任意の点列は
広義一様収束の意味で収束する部分列を持つか，
もしくは
広義一様収束の意味で
$\infty$に収束(普通はこれは発散と呼ぶけど)
する部分列を持つかの
いずれかが成り立つときにいう．
ジェネトポ風にいうと
$f\in \mathcal{F}$の定義域
を$f\colon D\to \cc\hookrightarrow\hat{\cc}$
というふうにリーマン球面$\hat{\cc}$
に拡張したと考えた時に
$\mathcal{F}$がコンパクト開位相における
相対コンパクト集合になるときにいうということである．
ただし有理型関数に収束することを許さず
そういうときには無限大という定数関数に収束すると約束している点に注意しよう．
有理型関数にも同様の定義ができてそれは
もうコンパクト開位相で相対コンパクトで良い．
\textbf{ブロッホの原理}
とは複素関数に関する性質$P$に関して以下の二つが同値であることを主張している
\textbf{指導原理}である．
\begin{enumerate}
\item 
大域関数$f\colon \cc\to \cc$で性質$P$
を満たすものは定数に限る．
\item 
領域$D\subset \cc$に関して
族$\{\, f\mid \text{$f$は$D$上の複素関数で$P$を満たす}\, \}$
は正規である．
\end{enumerate}
\textbf{この二つの同値性は一般には正しくない．}
しかし
いい感じの性質に関してはちゃんと成り立っており，
何やら裏がありそうという訳で研究されているのである．
先ほどから引用している
\cite{MR0379852}
では性質$P$が上の同値性を満たすための十分条件を与えており，おそらく最近ではZalcmanの原理の方が有名なんじゃないんだろうか？そこら辺の機微はよくわからない．
詳細は省略する．
ここら辺の話はとても面白いと思う．




\subsection{シュワルツの補題}
細かいことは
大体
\cite{MR1704258}
に書いてあると思う．
解析関数の剛性に関する話題はかなりたくさんあると思うが，
シュワルツの補題もその一つである．
以下の主張がそうである．
\begin{thm}
$r_{1}, r_{2}\in (0, \infty)$
とし
$f$を$\{|z|<R_{1}\}$
上で定義された正則関数とする．
もしも$f(0)=0$であり，
定義域上で$|f(z)|<r_{2}$を満たすならば
\[
|f(z)|\le \frac{r_{2}}{r_{1}}|z|
\]
が成り立つ．
さらにもしもある0でない点$z$
で
\[
|f(z)|=\frac{r_{2}}{r_{1}}|z|
\]
が満たされるのであれば，
ある実数$\alpha\in \rr$
が存在して任意の$z$で
\[
f(z)=\frac{r_{2}}{r_{1}}\exp(\sqrt{-1}\alpha)z
\]
が成り立つ．
\end{thm}
シュワルツ・ピックの定理というのは
これの$f(0)=0$
の仮定を外して開円盤の双正則同型を同定する話．


このシュワルツの補題を
リーマン計量の話に拡張したのが
アールフォルス
\cite{MR1501949}
である．
いろいろ同値な言い換えがあるが，
$\lambda$を単位開円盤上のポアンカレ計量
とする．これは
\[
\lambda(z)=\frac{2}{(1-|z|^{2})}
\]
と書かれる
(本当は$\lambda(z)^{2}dz\otimes dz$
がポアンカレ計量と呼ばれると思うが)．
このとき$\lambda$
のガウス曲率は全ての点で$-1$になる．
つまりはこの計量を備えた開円盤というのは
(二次元)双曲空間の円盤モデルと呼ばれるものに等しいということである．
アールフォルスの補題は
$|z|<1$上の退化しててもいいリーマン計量$\rho(z)$
のガウス曲率が$-1$
より小さいと
必ず$|\rho|\le \lambda$
を満たすということを言っている．
なぜだか知らないがこれを使うと色々なことがわかってすごいらしい．
論文\cite{MR0725266}
も参照のこと．
他には論文
\cite{MR0690493}, 
\cite{MR0486659}, 
\cite{MR2492498}, 
\cite{MR2120588}, 
\cite{MR1694951}, 
\cite{MR0908253}, 
\cite{MR3451773}, 
および
\cite{MR2745751}
を参照のこと．
はてなブログ
\cite{hahaha1}
も参照のこと．



\subsection{ピカールの定理}
ノート
\cite{pipi2015}
にいろいろ書いてある．


\begin{thm}
ピカールの定理は二つある．
\begin{enumerate}
\item (ピカールの小定理)：
大域関数$f\colon \cc\to \hat{\cc}$に関して
もしも$f$が定数でないのならば
$\hat{\cc}\setminus f(\cc)$
の濃度は高々$2$になる．
つまり$f$は高々$2$点を除いて全ての(無限大を含めた)複素数値を取りうる．

\item (ピカールの大定理)：
正則関数$f$について$z_{0}$
は真性特異点とする．
このとき$f$は$z_{0}$
の近傍で高々$1$つの値を除いて
全ての複素数値を無限回取る
($f^{-1}(w)$が無限集合って意味)．

\end{enumerate}
\end{thm}

昔の「日本数学物理学会誌」
では外国の論文を翻訳していたようで
\cite{192746}
はランダウによるピカールの定理の証明らしい．

論文
\cite{MR1546112}
ではランダウの定理とかピカールの定理などを統一するような定理を証明しているようだ．すごい．


\subsection{ショットキーの定理}
ショットキーの定理というのは
$f=a_{0}+a_{1}z+\dots$
が$|z|<1$で正則で
$0$と$1$の値を取らないならば
$|a_{1}|$が上から$a_{0}$に依存する定数で抑えられるという定理である．
以下の論文がそういう論文．
\cite{MR0021590}, 
\cite{MR0066460}, 
\cite{MR0575385}. 
論文
\cite{MR1574949}
もそういうような話題．
これの解説のようなものが
\cite{1927224}
にある．


\subsection{複素関数基礎編}
上で述べた
ランダウ，モンテル，
ピカール，ショットキーの定理に統一的視点を与えている論文を紹介する．

大体
\cite{MR0725266}
を読むといいと思う．
論文
\cite{MR0725266}
では
$\cc_{0, 1}=\hat{\cc}\setminus \{0, 1, \infty\}$
に対してその上の計量を
\[
\alpha(w)=\frac{(1+|w|^{1/3})^{1/2}}{|w|^{5/6}}
\times \frac{(1+|w-1|^{1/3})^{1/2}}{|w-1|^{5/6}}
\]
と定めるとある定数$k>0$
があって
$\alpha$の曲率は$-k$
以下になることをゴリゴリ計算して証明して
$\tau=\sqrt{k}\alpha$
とおくことで
$\cc_{0, 1}$上の計量$\tau$
でその曲率が$-1$以下になるものを
構成している．
そしてこの$\cc_{0, 1}$
上の計量とアールフォルスの定理を用いると
なぜだか知らないが
ランダウの定理やらピカールやらが証明できる．
すごい．

論文
\cite{MR0718986}
もショットキーの定理やらなんやらの古典的な
定理のある意味での同値性を与えていて
興味深い．

\section{その他}
なんか余ったやつ．

\subsection{Valironの定理}
これはよくわかんなかった．
多分複素力学系の話だと思う．
\cite{MR2058109}. 


\subsection{複素関数論の実証明}
論文\cite{MR0328028}
は複素関数論の定理を実解析で証明する話だと思う．
論文
\cite{MR2827550}
にも注目．


\subsection{オスグッド}
オスグッドの定理はたくさんあるらしい．

論文
\cite{MR1502274}
では一変数の実解析関数の無限和が解析的みたいな話をしている．

論文
\cite{MR0220914}
では領域上の正則関数の境界の拡張の話，つまりカラテオドリの定理もオスグッドが独立にやっていたらしいことが紹介されている．

論文
\cite{MR0042500}
では
各点収束と連続性の話，
なんか実連続関数の列の各点収束先が連続だと部分的に
広義一様収束が言えるみたいなことが書いてある．

論文
\cite{MR1500577}
では変分問題に手を出していたらしいことが書いてある．


論文
\cite{MR2794615}
は他変数解析関数の話で
オスグッドとかハルトークスとかの話．

\section*{参考文献たち}



\renewcommand{\refname}{ブログ記事}
	
\begin{thebibliography}{99}

\bibitem[はてな1]{hahaha1}
はてなブログ「おぼえがき」の記事「Ahlforsの観たSchwarzの補題」

https://hackberryfield.hatenablog.com/entry/2018/12/17/154219
\end{thebibliography}


\renewcommand{\refname}{一般の参考文献}
%READ!
%myplain is the same to amsplain style. 
\nocite{*}
\bibliographystyle{myplaindoidoi}
\bibliography{complexa.bib}



\end{document}